\documentclass[11pt]{article}

\usepackage[T1]{fontenc}
\usepackage[utf8]{inputenc}
\usepackage[margin=1in]{geometry}
\usepackage{booktabs}
\usepackage{hyperref}
\usepackage{longtable}
\usepackage{enumitem}
\usepackage{xcolor}

\hypersetup{
  colorlinks=true,
  linkcolor=blue,
  citecolor=blue,
  urlcolor=blue
}
\urlstyle{same}
\setlist[itemize]{leftmargin=*}
\setlist[enumerate]{leftmargin=*}

\title{MisinfoEQA: Stress-Testing Evidence Quality and Evaluation Fragility in Misinformation Benchmarks}
\author{
Fortune Nwachukwu Onwe\\
MSc Artificial Intelligence and Data Science\\
University of Hull, United Kingdom\\
\href{mailto:f.onwe-2025@hull.ac.uk}{f.onwe-2025@hull.ac.uk}\\
\href{mailto:fortonwe@gmail.com}{fortonwe@gmail.com}\\
\href{https://www.linkedin.com/in/fortune-onwe/}{linkedin.com/in/fortune-onwe}\\
\href{https://github.com/FortOnwe}{github.com/FortOnwe}\\
\href{https://github.com/FortOnwe/misinfo-eqa}{github.com/FortOnwe/misinfo-eqa}
}
\date{April 22, 2026}

\begin{document}
\maketitle

\begin{abstract}
Misinformation datasets are often evaluated as if labels, evidence, and splits
are directly comparable across sources. In practice, evidence fields may contain
natural-language rationales, structured references, long fact-checking articles,
or no usable evidence at all. We present MisinfoEQA, a lightweight evaluation
quality assurance harness that normalizes misinformation datasets, separates raw
evidence availability from natural-language evidence usability, runs simple
interpretable baselines, applies stress tests, and produces manual-audit sheets.
On four subsets from \texttt{ComplexDataLab/Misinfo\_Datasets} with 1,000
examples per dataset, MisinfoEQA finds that Climate-FEVER has high
natural-language evidence coverage but claim-plus-evidence underperforms
claim-only by 0.109 macro-F1, and manual audit confirms 76.0\% precision among
reviewed Climate-FEVER evidence flags. FEVER has 88.2\% raw evidence coverage
but 0.0\% natural-language evidence coverage in this normalized source, because
the evidence is mostly structured reference metadata. PubHealthTab has no usable
evidence in the sampled records. Snopes has complete natural-language evidence
coverage but low audit precision for simple evidence heuristics, showing that
long-form fact-checking narratives need richer semantic or article-structure-aware
QA. These results support a simple conclusion: evidence availability is not
evidence usability, and misinformation benchmark reports should include
dataset-QA diagnostics alongside model scores.
\end{abstract}

\section{Introduction}

Recent work on misinformation evaluation has expanded from static fact-checking
classification to richer concerns such as dataset quality, LLM resilience,
temporal generalization, and contamination resistance. A broad guide to
misinformation detection data and evaluation curated a large unified collection
and highlighted quality concerns such as keyword shortcuts, temporal artifacts,
and feasibility of verification \cite{thibault2025guide,complexdatalab2025misinfo}.
Adjacent benchmark work such as MisinfoBench \cite{yang2025misinfobench} and
TripleFact \cite{xu2025triplefact} targets LLM resilience and contamination
risks.

MisinfoEQA is complementary. It does not introduce a new LLM benchmark or a new
misinformation generation task. Instead, it asks whether a dataset is safe to
evaluate in the way a practitioner might be tempted to evaluate it. Are labels
mapped consistently? Does the evidence field contain natural-language evidence
or only IDs? Does adding evidence improve or hurt a lightweight classifier? Do
model rankings change on ambiguous slices? Do heuristic rationale flags survive
manual audit?

The contribution is a standalone command-line harness and an MVP validation
study. The tool:

\begin{enumerate}
  \item normalizes public or local datasets into a shared schema;
  \item distinguishes raw evidence coverage from natural-language evidence coverage;
  \item runs dependency-light baselines and stressors;
  \item generates Markdown/HTML reports, plots, and JSON artifacts;
  \item creates audit sheets and summarizes reviewer verdicts.
\end{enumerate}

The empirical claim is deliberately modest but useful: simple QA checks surface
actionable differences between datasets before expensive model evaluation.

\section{Related Work}

FEVER introduced a large-scale claim verification benchmark with claims labeled
as supported, refuted, or not enough information, and evidence annotations for
supported and refuted claims \cite{thorne2018fever}. Climate-FEVER adapted
FEVER-style verification to real-world climate claims and emphasized the
subtlety of modeling climate claims with evidence \cite{diggelmann2020climatefever}.
The ComplexDataLab misinformation collection unifies many misinformation
datasets and harmonizes labels into true, false, and unknown categories
\cite{thibault2025guide,complexdatalab2025misinfo}.

Recent benchmark work has broadened the evaluation target. MisinfoBench
evaluates LLM resilience to misinformation across dimensions such as
discernment, resistance, and principled refusal \cite{yang2025misinfobench}.
TripleFact focuses on benchmark contamination and proposes dynamic and
controlled evaluation components for LLM-driven fake news detection
\cite{xu2025triplefact}. MisinfoEQA differs by placing the dataset itself under
test. It is designed to run before, or alongside, model benchmarking.

\section{System Design}

\subsection{Normalized Schema}

Every record is coerced into:

\begin{quote}
\texttt{id, dataset, split, claim, evidence\_text, label3, date, source\_url, metadata}
\end{quote}

\texttt{label3} is restricted to \texttt{true}, \texttt{false}, and
\texttt{unknown}. Dataset-specific fields are preserved in \texttt{metadata}
when they are not mapped into the common schema.

\subsection{Evidence Typing}

MisinfoEQA reports two evidence notions: raw evidence coverage, where any
evidence-like value is present, and natural-language evidence coverage, where
the value appears to be usable text. This distinction prevents structured tuples
such as FEVER evidence references from being treated as ordinary evidence text.

\subsection{Baselines}

The MVP uses lightweight baselines: a majority-class classifier; hashed TF-IDF
plus softmax logistic regression on claim text, evidence text, and claim plus
evidence; and a heuristic NLI-style cue classifier. The TF-IDF classifier is
dependency-free so the project can run in small environments and continuous
integration without a heavy ML stack.

\subsection{Stressors}

MisinfoEQA implements five stressors: keyword shortcut masking, temporal shift
when dates exist, evidence ablation, ambiguity slicing, and label-rationale
mismatch checks. The label-rationale check flags missing evidence, low local
evidence relevance, and guarded label-cue disagreements.

\subsection{Manual Audit}

The harness writes \texttt{audit\_sheet.csv} from flagged examples. Reviewers
mark each row as \texttt{real\_dataset\_issue}, \texttt{evidence\_mismatch},
\texttt{missing\_or\_weak\_evidence}, \texttt{structured\_reference\_evidence},
\texttt{heuristic\_false\_positive}, \texttt{baseline\_limitation}, or
\texttt{unclear}. Summary precision treats the first four labels as real issues.

\section{Experimental Setup}

The MVP evaluates four subsets from
\texttt{ComplexDataLab/Misinfo\_Datasets}: \texttt{fever},
\texttt{climate\_fever}, \texttt{pubhealthtab}, and \texttt{snopes}. The
canonical run uses 1,000 examples per dataset, seed 42, no API judge, and 1,000
bootstrap samples.

The command sequence is:

\begin{verbatim}
python -m misinfo_eqa scan --config configs/mvp.yaml
python -m misinfo_eqa run --config configs/mvp.yaml
python -m misinfo_eqa audit --latest --per-dataset 25
python -m misinfo_eqa audit-summary --latest
\end{verbatim}

The local validation run used for this paper was \texttt{runs/20260422-040519}.

\section{Results}

\subsection{Coverage}

\begin{table}[h]
\centering
\small
\begin{tabular}{lrrrr}
\toprule
Dataset & Examples & Natural-lang. evidence & Raw evidence & Date coverage \\
\midrule
\texttt{climate\_fever} & 1000 & 99.8\% & 100.0\% & 0.0\% \\
\texttt{fever} & 1000 & 0.0\% & 88.2\% & 0.0\% \\
\texttt{pubhealthtab} & 1000 & 0.0\% & 0.0\% & 0.0\% \\
\texttt{snopes} & 1000 & 100.0\% & 100.0\% & 0.0\% \\
\bottomrule
\end{tabular}
\caption{Coverage summary for the canonical MVP run.}
\label{tab:coverage}
\end{table}

The strongest coverage finding is that raw evidence coverage can be misleading.
FEVER has raw evidence in 88.2\% of sampled records, but those fields are
structured references rather than natural-language evidence.

\subsection{Baseline Metrics}

\begin{table}[h]
\centering
\small
\begin{tabular}{lrrrrr}
\toprule
Dataset & Majority & Claim & Claim+Evidence & Evidence & Heuristic \\
\midrule
\texttt{climate\_fever} & 0.228 & 0.660 & 0.551 & 0.407 & 0.207 \\
\texttt{fever} & 0.217 & 0.374 & 0.367 & 0.217 & 0.168 \\
\texttt{pubhealthtab} & 0.214 & 0.391 & 0.391 & 0.214 & 0.125 \\
\texttt{snopes} & 0.245 & 0.312 & 0.329 & 0.359 & 0.376 \\
\bottomrule
\end{tabular}
\caption{Macro-F1 by lightweight baseline.}
\label{tab:baselines}
\end{table}

The Climate-FEVER result is the clearest evidence-ablation signal: adding
evidence reduces macro-F1 from 0.660 to 0.551. This does not prove that evidence
is useless; rather, it shows that the evidence field and lightweight classifier
interact poorly enough to justify closer QA.

\subsection{Stressor Findings}

\begin{table}[h]
\centering
\small
\begin{tabular}{lrrrr}
\toprule
Dataset & Keyword drop & Evidence gain & Ambiguity Spearman & Label-rationale flags \\
\midrule
\texttt{climate\_fever} & 0.017 & -0.109 & 0.400 & 185/315 \\
\texttt{fever} & 0.007 & skipped & 0.000 & skipped \\
\texttt{pubhealthtab} & 0.004 & skipped & 0.000 & skipped \\
\texttt{snopes} & 0.013 & 0.016 & 0.900 & 75/301 \\
\bottomrule
\end{tabular}
\caption{Selected stressor outputs. Evidence gain is claim-plus-evidence macro-F1 minus claim-only macro-F1.}
\label{tab:stressors}
\end{table}

Temporal shift was skipped for all four datasets because date coverage was
0.0\% in the sampled normalized records.

\subsection{Manual Audit}

\begin{table}[h]
\centering
\small
\begin{tabular}{lrrrl}
\toprule
Dataset & Audited & Real issues & Precision & Verdict mix \\
\midrule
\texttt{climate\_fever} & 25 & 19 & 76.0\% & mismatch:8; weak:11; false positive:6 \\
\texttt{snopes} & 25 & 4 & 16.0\% & real:2; weak:2; false positive:21 \\
\bottomrule
\end{tabular}
\caption{Manual audit results for flagged examples.}
\label{tab:audit}
\end{table}

The manual audit validates Climate-FEVER as the strongest MVP finding. It also
reveals a limitation: the Snopes flags are mostly false positives because long
fact-checking narratives contain background, quoted claims, and verdict language
that simple lexical heuristics do not reliably isolate.

\section{Discussion}

MisinfoEQA changes the first question from ``Which model wins?'' to ``Is this
dataset being evaluated in a way that matches its fields?'' This is useful
because the four MVP datasets fail or stress different assumptions:

\begin{itemize}
  \item Climate-FEVER has usable-looking evidence, but evidence hurts a
  lightweight classifier and audited flags frequently identify weak or mismatched
  evidence.
  \item FEVER requires resolving structured evidence references into sentence
  text before evidence-based evaluation is meaningful.
  \item PubHealthTab cannot support evidence-ablation or rationale checks from
  this normalized source.
  \item Snopes contains natural-language evidence, but long-form articles need
  stronger span selection.
\end{itemize}

The results suggest that evaluation reports should include dataset-QA tables
before presenting model rankings. In particular, evidence coverage should be
reported as natural-language evidence coverage, not just non-empty evidence
fields.

\section{Limitations}

The MVP uses 1,000 sampled examples per dataset, not full-dataset sweeps. It uses
lightweight baselines, not modern neural or API judges. Temporal shift is skipped
because the sampled records have no usable dates. Manual audit covers 50
examples. The Snopes heuristic needs better article span extraction or semantic
validation. Results may shift if the upstream Hugging Face collection changes.

\section{Ethical Considerations}

MisinfoEQA is designed to inspect existing public datasets. It does not generate
new misinformation, retrieve live claims, or create adversarial false examples.
Reports may include snippets from existing datasets, so users should handle
outputs with the same care as the source data and avoid presenting flagged
claims without context.

\section{Conclusion}

MisinfoEQA provides a small, reproducible harness for stress-testing
misinformation dataset quality. The MVP shows that evidence fields vary sharply
in usability, that evidence can hurt simple evaluation setups, and that manual
audit is necessary to separate true dataset issues from heuristic artifacts. The
next version should resolve structured evidence references, add stronger
semantic rationale checks, expand manual audit, and compare light QA diagnostics
against stronger LLM or NLI judges.

\bibliographystyle{plain}
\bibliography{references}

\end{document}
